\hspace*{\fill} 

\cvsection{研究兴趣}
\cvevent{\textbf{计算机图形学}}{材质迁移、纹理迁移}{}{}

\cvsection{技能特长}
\cvevent{\textbf{图形学基础理论}}{}{}{}
\begin{itemize}
\item 曾担任深圳大学计算机图形学课程助教，对图形学基础知识以及OpenGL编程有一定的了解。
\end{itemize}

\divider

\cvevent{\textbf{深度学习实践}}{}{}{}
\begin{itemize}
\item 熟悉主流深度学习框架（Pytorch、Tensorflow等）,
\item 了解生成生成对抗网络，三元网络，图神经网络等多种类型的神经网络，并在原有的网络结构和损失函数上进行优化。
\end{itemize}

\divider

\cvevent{\textbf{其他}}{}{}{}
\begin{itemize}
\item 熟悉Python编程，
\item 熟悉三维建模软件的使用（Blender、Meshlab等），
\item 熟悉常用的图像处理与视频剪辑软件，
\item 参与论文投稿、专利申请等多种类型的文档写作。
\end{itemize}

\cvsection{校园经历}
\begin{itemize}
\item 深大计软学院研究生第二党支部宣传委员
\item 沈建大信息学院计算机1502班班长
\item 沈建大信息学院计算机1702班助理班主任
\item 沈阳市大学生记者团新媒体工作部干事
\end{itemize}


% \cvsection{校园经历}
% \cvtag{沈阳建筑大学计算机1502班班长} 
% \cvtag{沈阳建筑大学计算机1702班助理班主任}
% \cvtag{沈阳建筑大学信息学院社管部部长}\newline
% \cvtag{沈阳市大学生记者团新媒体工作部记者}

% \cvsection{所获荣誉}
% \cvtag{沈阳建筑大学优秀毕业生} 
% \cvtag{沈阳建筑大学校友联系大使}
% \cvtag{沈阳建筑大学优秀学生干部}
% \cvtag{沈阳建筑大学社团工作先进个人}






\begin{comment}
\cvsection{Languages}
 
\textcolor{accent}  \faLanguage \hspace{} \hspace{} \textcolor{emphasis}{\textbf{ English}}

\divider

\cvskill{\textcolor{accent}  \faLanguage \hspace{} \hspace{} Hindi}{5}
\divider

\cvskill{\textcolor{accent}  \faLanguage \hspace{} \hspace{} German}{3}

%% Yeah I didn't spend too much time making all the 
%% spacing consistent... sorry. Use \smallskip, \medskip, 
%% \bigskip, \vpsace etc to make ajustments.
\medskip

\cvsection{Education}

\cvevent{MBA\ in Finance \& Marketing}{IMS, DAVV, Indore}{Jun 2005 -- June 2007}{}


\divider

\cvevent{MA\ Economics}{University of Pune}{June 2004 -- June 2006}{}

\divider
\end{comment}